% Generated draft: requires mathematical and source review.
\documentclass[11pt]{article}
\usepackage[margin=25mm]{geometry}
\usepackage{fontspec}
\setmainfont{Latin Modern Roman}
\usepackage{amsmath,amssymb,mathtools,unicode-math}
\setmathfont{Latin Modern Math}
\usepackage{enumitem,needspace,xurl,hyperref}
\setlength{\parindent}{0pt}
\setlength{\parskip}{5pt}
\setlength{\emergencystretch}{3em}
\newcommand{\N}{\mathbb{N}}
\newcommand{\Z}{\mathbb{Z}}
\newcommand{\Q}{\mathbb{Q}}
\newcommand{\R}{\mathbb{R}}
\newcommand{\C}{\mathbb{C}}
\newcommand{\F}{\mathbb{F}}
\newcommand{\A}{\mathbb{A}}
\newcommand{\PP}{\mathbb{P}}
\newcommand{\E}{\mathbb{E}}
\newcommand{\eps}{\varepsilon}
\newcommand{\dd}{\,\mathrm{d}}
\begin{document}
\section*{UnsolvedMath ID: 6700073}
\textbf{Source problem code:} catalogue entry
\subsection{Definitions and mathematical statement}
Fix a \(C^2\)-smooth Riemannian manifold \((W,g)\) and a closed subset \(Y\subset W\). Characterize when there is a family of smooth domains \(V_i\subset W\), indexed by an ordered or directed set \(I\), such that
\[
 Y=\bigcap_{i\in I}V_i,
 \qquad \operatorname{Scal}(\partial V_i)\longrightarrow+\infty.
\]
The scalar curvature is that of the hypersurface with its metric induced from \(g\), rather than ambient scalar curvature or mean curvature. Under a uniform interpretation of blow-up, the second condition means
\[
 \forall A>0\ \exists i_0\ \forall i\succeq i_0\ \forall p\in\partial V_i:
 \quad\operatorname{Scal}_{g|_{\partial V_i}}(p)\ge A.
\]
The excerpt does not specify uniform versus locally uniform blow-up on noncompact boundaries, so that convention needs confirmation.

For noncompact \(Y\), the source also proposes stronger versions: one may require the domains to be decreasing,
\[
 j\succeq i\quad\Longrightarrow\quad V_j\subseteq V_i,
\]
or additionally require eventual containment in every neighborhood,
\[
 \forall U\supseteq Y\text{ open}\ \exists i_0\ \forall i\succeq i_0:
 \quad V_i\subseteq U.
\]
These strengthen the intersection requirement and are not declared equivalent to it. The source says the latter neighborhood requirement calls for uncountable indexing, so the problem is not silently restricted to sequences. The exact index order, open-versus-closed domain convention, permitted domain compactness, and empty-boundary convention are not given.

The passage announces subsequent conjectures but ends with an incomplete reference marker. The supplied summary describes expected low-Hausdorff-dimension examples but gives no dimensional threshold or complete conjecture; none is invented as a characterization. It reports the general ``remnants of collapse'' program open with no resolution located. This is source information without independently verified current status.
\subsection{Short English statement}
Which closed subsets can be intersections of smooth domains whose induced boundary scalar curvatures tend to positive infinity? The source also asks about nested and neighborhood-cofinal versions for noncompact sets, allowing uncountable indexing; the precise blow-up and domain conventions need review.
\subsection{Sources}
\begin{itemize}
\item Local registry: \url{data/unsolvedmath/problems.json}, permanent UnsolvedMath ID \(6700073\), source code catalogue entry; supplied statement, background, and research summary, updated 2026-08-06.
\item Misha Gromov, \emph{101 Questions, Problems and Conjectures around Scalar Curvature} (2017), printed tag [?77], occurrence 76, PDF page 77: \url{https://www.ihes.fr/~gromov/wp-content/uploads/2018/08/101-problemsOct1-2017.pdf}. Supplied normalized PDF-text extraction accessed 2026-07-29; described as incomplete.
\item The announced following conjectures are not included in full, and no separate resolution reference is supplied.
\end{itemize}
\end{document}
